% Options for packages loaded elsewhere
\PassOptionsToPackage{unicode}{hyperref}
\PassOptionsToPackage{hyphens}{url}
\PassOptionsToPackage{dvipsnames,svgnames,x11names}{xcolor}
\documentclass[
  a4paper,
]{article}
\usepackage{xcolor}
\usepackage[margin=25mm]{geometry}
\usepackage{amsmath,amssymb}
\setcounter{secnumdepth}{-\maxdimen} % remove section numbering
\usepackage{iftex}
\ifPDFTeX
  \usepackage[T1]{fontenc}
  \usepackage[utf8]{inputenc}
  \usepackage{textcomp} % provide euro and other symbols
\else % if luatex or xetex
  \usepackage{unicode-math} % this also loads fontspec
  \defaultfontfeatures{Scale=MatchLowercase}
  \defaultfontfeatures[\rmfamily]{Ligatures=TeX,Scale=1}
\fi
\usepackage{lmodern}
\ifPDFTeX\else
  % xetex/luatex font selection
  \setmainfont[]{DejaVu Serif}
  \setsansfont[]{DejaVu Sans}
  \setmonofont[]{DejaVu Sans Mono}
\fi
% Use upquote if available, for straight quotes in verbatim environments
\IfFileExists{upquote.sty}{\usepackage{upquote}}{}
\IfFileExists{microtype.sty}{% use microtype if available
  \usepackage[]{microtype}
  \UseMicrotypeSet[protrusion]{basicmath} % disable protrusion for tt fonts
}{}
\makeatletter
\@ifundefined{KOMAClassName}{% if non-KOMA class
  \IfFileExists{parskip.sty}{%
    \usepackage{parskip}
  }{% else
    \setlength{\parindent}{0pt}
    \setlength{\parskip}{6pt plus 2pt minus 1pt}}
}{% if KOMA class
  \KOMAoptions{parskip=half}}
\makeatother
\usepackage{longtable,booktabs,array}
\newcounter{none} % for unnumbered tables
\usepackage{calc} % for calculating minipage widths
% Correct order of tables after \paragraph or \subparagraph
\usepackage{etoolbox}
\makeatletter
\patchcmd\longtable{\par}{\if@noskipsec\mbox{}\fi\par}{}{}
\makeatother
% Allow footnotes in longtable head/foot
\IfFileExists{footnotehyper.sty}{\usepackage{footnotehyper}}{\usepackage{footnote}}
\makesavenoteenv{longtable}
\setlength{\emergencystretch}{3em} % prevent overfull lines
\providecommand{\tightlist}{%
  \setlength{\itemsep}{0pt}\setlength{\parskip}{0pt}}
\usepackage{luatexja}
\usepackage{xurl}
\usepackage{bookmark}
\IfFileExists{xurl.sty}{\usepackage{xurl}}{} % add URL line breaks if available
\urlstyle{same}
\hypersetup{
  pdftitle={Emergence of Spacetime and Elementary Particles from a Bounded Hilbert Space and Octonions},
  pdfauthor={Noriaki Kihara (WF System Co., Ltd.) ORCID 0009-0004-6753-4020},
  colorlinks=true,
  linkcolor={Maroon},
  filecolor={Maroon},
  citecolor={Blue},
  urlcolor={Blue},
  pdfcreator={LaTeX via pandoc}}

\title{Emergence of Spacetime and Elementary Particles from a Bounded
Hilbert Space and Octonions}
\usepackage{etoolbox}
\makeatletter
\providecommand{\subtitle}[1]{% add subtitle to \maketitle
  \apptocmd{\@title}{\par {\large #1 \par}}{}{}
}
\makeatother
\subtitle{― From Photon Geometry and Gnomonic Central Projection to Open
Problems in Background-Free Quaternionic and Octonionic Dynamics ―}
\author{Noriaki Kihara (WF System Co., Ltd.) ORCID 0009-0004-6753-4020}
\date{November 16, 2026 Version DOI 10.5281/zenodo.24267348}

\ltjsetparameter{jacharrange={-2,-3}}
\begin{document}
\maketitle

\textbf{Author:} Noriaki Kihara (WF System Co., Ltd.) ORCID
0009-0004-6753-4020\\
\textbf{Date:} 2026-11-16 \textbf{Version:} v1.0\\
\textbf{Version DOI:} 10.5281/zenodo.24267348 \textbf{Concept DOI:}
10.5281/zenodo.24267347\\
\textbf{License:} CC BY 4.0

\begin{center}\rule{0.5\linewidth}{0.5pt}\end{center}

\subsection{Abstract}\label{abstract}

This manuscript is a record of thought experiments and dialogue starting
from a geometric intuition about photons (being stationary in
four-dimensional spacetime) and proceeding through a central-projection
model, retention of phase gradients by quaternions, and the
non-associativity of octonions toward the ``self-organization of
particles and interactions from a Hilbert space without assuming an
external background spacetime.''

\begin{center}\rule{0.5\linewidth}{0.5pt}\end{center}

\subsection{Introduction: Positioning of This Manuscript and Connection
to Existing
Research}\label{introduction-positioning-of-this-manuscript-and-connection-to-existing-research}

This manuscript records a thought experiment that attempts to
reconstruct spacetime, particles, and interactions solely from state
relations in a bounded, finite Hilbert space and the algebraic
structures of quaternions and octonions, without positing a
four-dimensional background spacetime as a stage.

At first glance, such a starting point appears far removed from standard
textbook formulations; however, both its problem consciousness and its
mathematical apparatus connect to several active currents in
contemporary theoretical physics:

\begin{enumerate}
\def\labelenumi{\arabic{enumi}.}
\tightlist
\item
  \textbf{Emergent Spacetime and Bounded State Spaces:} It shares a
  common problem consciousness with approaches that reconstruct spatial
  dimension and geometry solely from correlations among quantum states
  and decompositions into subsystems without assuming a 4-dimensional
  manifold (such as Quantum Mereology and tensor-network emergent
  spacetime), as well as discussions examining finite-dimensional,
  bounded Hilbert spaces in correspondence with the event horizon of de
  Sitter spacetime.
\item
  \textbf{Spinors and Division Algebras for Spacetime and Elementary
  Particles:} Viewing spatial coordinates and light rays through
  complex-number pairs (spinors) and the Hopf fibration {[}X5{]}{[}X6{]}
  shares its mathematical structure with twistor theory and the
  Bloch-sphere representation in quantum information, while using
  octonionic non-associativity, triality, and minimal left ideals to
  describe particle species and symmetries connects to algebraic
  approaches to the Standard Model via division algebras
  {[}X8{]}{[}X9{]}.
\end{enumerate}

Note that this manuscript is not a claim of a completed theory, but an
ongoing organization of thought. The boundary between established
mathematical identities/physical interpretations, the working hypotheses
specific to this manuscript, and open derivation problems (the concrete
update rule that rewrites internal phases {[}S9{]}--{[}S11{]} and the
mathematical proof of factorization from \(N^2\) all-to-all coupling) is
strictly separated as shown in the ``Classification'' table at the end.

{\def\LTcaptype{none} % do not increment counter
\begin{longtable}[]{@{}
  >{\raggedright\arraybackslash}p{(\linewidth - 6\tabcolsep) * \real{0.2500}}
  >{\raggedright\arraybackslash}p{(\linewidth - 6\tabcolsep) * \real{0.2500}}
  >{\raggedright\arraybackslash}p{(\linewidth - 6\tabcolsep) * \real{0.2500}}
  >{\raggedright\arraybackslash}p{(\linewidth - 6\tabcolsep) * \real{0.2500}}@{}}
\toprule\noalign{}
\begin{minipage}[b]{\linewidth}\raggedright
Stage
\end{minipage} & \begin{minipage}[b]{\linewidth}\raggedright
Mathematical Structure / Number System
\end{minipage} & \begin{minipage}[b]{\linewidth}\raggedright
Geometric Map
\end{minipage} & \begin{minipage}[b]{\linewidth}\raggedright
Physical Correspondence (Working Hypothesis)
\end{minipage} \\
\midrule\noalign{}
\endhead
\bottomrule\noalign{}
\endlastfoot
Stage 1 & Minkowski metric \((x,y,z,t,\tau)\) & Cross-section of 5D
light cone & Zero proper time for photons (\(ds=0\)); temporal
progression of matter \\
Stage 2 & Imaginary time \(R^2 = t^2 + (i\tau)^2\) & Central projection
(Beltrami coordinates) & de Sitter spacetime; event horizon \(R\) and
gravitational distortion \\
Stage 3 & Quaternions \(\mathbb{H}\) (conjugate pair of complex numbers)
& Hopf fibration \(S^1 \hookrightarrow S^3 \to S^2\) & Retention of 3D
phase gradients without external coordinates (photons) \\
Stage 4 & Octonions \(\mathbb{O}\) (conjugate pair of quaternions) &
Hopf fibration \(S^3 \hookrightarrow S^7 \to S^4\) & Closed 4D spacetime
\(S^4\), 3 internal degrees of freedom, and self-organization of
condensates \\
\end{longtable}
}

\begin{center}\rule{0.5\linewidth}{0.5pt}\end{center}

\subsection{Stage 1: Photons and the Geometry of Four-Dimensional
Spacetime}\label{stage-1-photons-and-the-geometry-of-four-dimensional-spacetime}

\textbf{{[}Kihara{]}} Why do photons appear to always travel at the
speed of light?

\textbf{{[}Gemini{]}} The reason photons always travel at the speed of
light (approximately 300,000 km/s) is that their rest mass is zero. From
different perspectives in physics, this is organized as follows:

\begin{enumerate}
\def\labelenumi{\arabic{enumi}.}
\tightlist
\item
  \textbf{Perspective of Relativity:} In the relation between energy
  \(E\), momentum \(p\), and rest mass \(m_0\),
  \(E^2 = (pc)^2 + (m_0 c^2)^2\), massive matter (\(m_0 > 0\)) retains
  energy \(E = m_0 c^2\) even at rest and cannot reach \(c\) because the
  required energy diverges. For a photon (\(m_0 = 0\)), the relation
  reduces to \(E = pc\); to carry nonzero energy, a photon must always
  have momentum \(p\) and move at speed \(c\) (if it stopped,
  \(E = 0\)).
\item
  \textbf{Perspective of Particle Physics:} Particles that interact with
  the Higgs field (such as electrons and quarks) acquire mass, whereas
  photons do not interact with the Higgs field and propagate at the
  upper speed limit of causality, \(c\).
\item
  \textbf{Perspective of Electromagnetism:} The propagation speed of
  linked electric and magnetic fields (electromagnetic waves) in vacuum
  is determined solely by the vacuum permittivity \(\varepsilon_0\) and
  permeability \(\mu_0\) as \(c = 1/\sqrt{\varepsilon_0 \mu_0}\).
\item
  \textbf{Note (Propagation in Media):} Even when light appears slower
  in water or glass, each photon still travels at \(c\) in the vacuum
  between molecules; the apparent speed decreases due to repeated
  electromagnetic interactions (absorption/re-emission and interference)
  with electrons in the medium.
\end{enumerate}

\textbf{{[}Kihara{]}} May we interpret this as photons alone being
stationary in four-dimensional spacetime, while only protons and
electrons move at the speed of time?

\textbf{{[}Gemini{]}} That viewpoint corresponds to describing special
relativity in terms of ``four-dimensional spacetime geometry (Minkowski
spacetime).'' Whereas the previous explanation focused on how particles
behave in space (energy and mass), this describes the same phenomenon
from how they exist within the overall spacetime structure (geometry).

\begin{enumerate}
\def\labelenumi{\arabic{enumi}.}
\tightlist
\item
  \textbf{Protons and Electrons Move Along the Time Axis at \(c\):} In a
  geometric interpretation of 4-dimensional spacetime, every object can
  be viewed as moving through 4D spacetime at a constant speed \(c\). A
  spatially stationary proton or electron devotes 100\% of its motion to
  the time axis (speed \(c\)), and as spatial motion increases, its
  progression along the time axis is reduced (time dilation). Because a
  photon directs 100\% of its motion into spatial directions, its
  progression along the time axis is zero.
\item
  \textbf{Zero Spacetime Interval for Photons:} The spacetime interval
  \(ds\) in 4-dimensional spacetime is
  \[ds^2 = -(c\,dt)^2 + dx^2 + dy^2 + dz^2.\] During a time interval
  \(dt\), a photon travels a spatial distance
  \(\sqrt{dx^2 + dy^2 + dz^2} = c\,dt\), so \(ds = 0\) always holds and
  no proper time elapses along its worldline.
\item
  \textbf{Clarification of Expression:} Strictly speaking, a photon does
  not remain at a single fixed point in 4D spacetime coordinates; on a
  Minkowski diagram its coordinates \((x, t)\) change along the
  45-degree light cone, while the elapsed proper time (4-dimensional
  interval) along that trajectory is identically zero.
\end{enumerate}

\textbf{{[}Kihara{]}} That is a mapping from a light cone:
\(x^2 + y^2 + z^2 = t^2 + R^2\).

\textbf{{[}Gemini{]}} That equation (in natural units with \(c = 1\))
corresponds to viewing phenomena in 4-dimensional spacetime as a
``cross-section (mapping) from a higher-dimensional light cone (null
structure).''

\begin{enumerate}
\def\labelenumi{\arabic{enumi}.}
\tightlist
\item
  \textbf{Cross-Section (Mapping) of a Higher-Dimensional Light Cone:}
  The squared 3D spatial distance \(x^2 + y^2 + z^2\) balances the sum
  of squares of time \(t\) and an additional parameter \(R\), which can
  be interpreted as slicing a light cone in a 5-dimensional space at a
  specific value \(R\) and mapping it onto 4-dimensional spacetime.
\item
  \textbf{Invariant Hyperboloid:} Rearranging gives
  \(x^2 + y^2 + z^2 - t^2 = R^2\): \(R = 0\) yields the light cone
  itself, while \(R > 0\) yields a one-sheeted hyperboloid attached
  outside the light cone.
\item
  \textbf{Note (Correspondence with the ``Speed of Time''):} When
  expressing the intuition that massive matter progresses along time via
  the Pythagorean theorem, one uses the matter's proper time \(\tau\) as
  \(t^2 = (x^2 + y^2 + z^2) + \tau^2\). In this right triangle (Epstein
  geometry) with coordinate time \(t\) as the hypotenuse, spatial
  distance as the base, and proper time \(\tau\) as the height, photons
  (\(\tau = 0\)) lie on the light-cone boundary, while protons and
  electrons (\(\tau > 0\)) lie inside it.
\end{enumerate}

\textbf{{[}Kihara{]}} The expression using proper time \(\tau\),
\(x^2 + y^2 + z^2 = t^2 - \tau^2\), seems easier to match with standard
theory.

\textbf{{[}Gemini{]}} \((x^2 + y^2 + z^2) + \tau^2 = t^2\) corresponds
directly to the standard metric of special relativity, connecting the
light-cone mapping and proper time \(\tau\).

\begin{enumerate}
\def\labelenumi{\arabic{enumi}.}
\tightlist
\item
  \textbf{Relation with Coordinate Time \(t\) as the Hypotenuse:}
  Photons (\(\tau = 0\)) satisfy \(x^2 + y^2 + z^2 = t^2\) on the
  boundary of the light cone, while protons and electrons (\(\tau > 0\))
  have \(\tau^2\) subtracted from \(t^2\) in the squared spatial
  distance, confining them inside the light cone (the geometric reason
  massive particles cannot exceed \(c\)).
\item
  \textbf{Interpretation as a Mapping from a 5D Light Cone:}
  \((x^2 + y^2 + z^2) + \tau^2 = t^2\) has the exact form of a pure
  light cone in a 5-dimensional space with 3 spatial coordinates
  \((x, y, z)\) plus a fourth spatial-like component \(\tau\).
  Everything runs on this 5D light cone at \(c\); photons direct their
  velocity vector entirely into \((x, y, z)\), whereas massive particles
  direct part or all of it into \(\tau\) and appear projected onto 3D
  space.
\item
  \textbf{Connection to Quantum Mechanics:} In quantum field theory
  (properties of the Dirac equation), even massive particles exhibit
  luminal zigzag motion (Zitterbewegung) below the Compton wavelength
  scale, appearing slower than \(c\) only on average.
\end{enumerate}

\begin{center}\rule{0.5\linewidth}{0.5pt}\end{center}

\subsection{Stage 2: Central Projection and de Sitter
Spacetime}\label{stage-2-central-projection-and-de-sitter-spacetime}

\textbf{{[}Kihara{]}} If we express \(R^2 = t^2 + (i\tau)^2\), it
becomes a central projection \[x^2 + y^2 + z^2 = R^2,\] which looks like
a case with gravitational distortion governed by a local radius of
curvature \(R\).

\textbf{{[}Gemini{]}} The transformation
\[t^2 - \tau^2 = t^2 + (i\tau)^2 = R^2 \implies x^2 + y^2 + z^2 = R^2\]
corresponds to a geometric procedure bridging special relativity (flat
spacetime) and general relativity (curved spacetime).

\begin{enumerate}
\def\labelenumi{\arabic{enumi}.}
\tightlist
\item
  \textbf{Transition from Hyperbolic to Spherical Geometry via Imaginary
  Time \((i\tau)\):} Multiplying proper time by the imaginary unit \(i\)
  gives \((i\tau)^2 = -\tau^2\) (Wick rotation / imaginary time),
  converting the hyperbolic relation into the equation of a sphere of
  radius \(R\), \(x^2 + y^2 + z^2 = R^2\).
\item
  \textbf{Central Projection and Gravity via Local Curvature Radius
  \(R\):} Viewing \(x^2 + y^2 + z^2 = R^2\) as a central projection from
  a higher-dimensional space (gnomonic projection / Beltrami--Klein
  model {[}X1{]}) maps geodesics on the sphere to straight lines in the
  projected flat space {[}S1{]}. Even in an arbitrary gravitational
  field, cutting out a local inertial frame corresponds to approximating
  it by central projection onto the tangent space of a sphere with local
  radius of curvature \(R\).
\item
  \textbf{Contrast Between Photons and Matter:} From
  \(R^2 = t^2 + (i\tau)^2\), photons (\(\tau = 0\)) have \(R = t\) and
  travel on the horizon/light-cone boundary. For massive matter
  (\(\tau > 0\)), \((i\tau)^2 = -\tau^2\) shrinks the effective radius
  \(R\) relative to \(t\) (\(R < t\)), giving the interpretation that
  possessing proper time \(\tau\) locally contracts the curvature radius
  \(R\) and generates gravitational distortion.
\end{enumerate}

\textbf{{[}Kihara{]}} This is a model I conceived in April this year,
which makes it possible to understand de Sitter spacetime
{[}X2{]}---which can appear anisotropic---as a simple central projection
{[}S1{]}{[}S2{]}.

\textbf{{[}Gemini{]}} Reinterpreting de Sitter spacetime {[}X2{]}
(accelerating expansion driven by a cosmological constant) not in FLRW
coordinates but via central projection (projective / Beltrami
coordinates) from a higher-dimensional hypersurface clarifies the
following structure {[}S1{]}{[}S2{]}:

\begin{enumerate}
\def\labelenumi{\arabic{enumi}.}
\tightlist
\item
  \textbf{Straightening of Geodesics:} All geodesics in de Sitter
  spacetime (light rays and inertial particle trajectories) map to
  straight lines in the projected space.
\item
  \textbf{Visualization of the Event Horizon:} The cosmological event
  horizon (Hubble radius), representing the causal limit, appears
  directly as the spherical boundary \(x^2 + y^2 + z^2 = R^2\).
\item
  \textbf{Reduction of Apparent Non-Stationarity and Anisotropy:}
  Differences in spatial stretching or time passage depending on the
  observer's viewpoint reduce to which angle or cross-section of a
  symmetric 5D hyperboloid (or 4D sphere under imaginary time \(i\tau\))
  is being centrally projected.
\end{enumerate}

\begin{center}\rule{0.5\linewidth}{0.5pt}\end{center}

\subsection{Stage 3: Extension from Complex Numbers to Quaternions and
Octonions}\label{stage-3-extension-from-complex-numbers-to-quaternions-and-octonions}

\textbf{{[}Kihara{]}} In this model, if we represent photons not with
complex waves but with quaternions, doesn't it become an
easy-to-understand mapping?

\textbf{{[}Gemini{]}} Representing photons with quaternions
(\(\mathbb{H}\)) rather than complex waves (\(\mathbb{C}\)) in the
central-projection model unifies ``spacetime geometry (central
projection)'' and ``wave oscillation (phase and polarization)'' into a
single mapping.

\begin{enumerate}
\def\labelenumi{\arabic{enumi}.}
\tightlist
\item
  \textbf{Absorption of the 3 Spatial Axes into
  \((\mathbf{i}, \mathbf{j}, \mathbf{k})\):} A standard complex wave
  \(e^{i(\mathbf{k}\cdot\mathbf{x}-\omega t)}\) has only one imaginary
  unit \(i\), requiring external vectors for propagation direction and
  polarization. Using a quaternion
  \(q = t + x\mathbf{i} + y\mathbf{j} + z\mathbf{k}\)
  (\(\mathbf{i}^2=\mathbf{j}^2=\mathbf{k}^2=\mathbf{i}\mathbf{j}\mathbf{k}=-1\)),
  the spatial vector
  \(\mathbf{r} = x\mathbf{i} + y\mathbf{j} + z\mathbf{k}\) satisfies
  \(\mathbf{r}^2 = -(x^2 + y^2 + z^2)\), so the sign flip of the
  Minkowski metric
  \(t^2 + \mathbf{r}^2 = t^2 - (x^2 + y^2 + z^2) = \tau^2\) is built
  into the algebraic structure.
\item
  \textbf{Rotation Mapping on \(S^3\):} By Euler's formula for a unit
  pure-imaginary quaternion \(\mathbf{n}\) (\(\mathbf{n}^2 = -1\)),
  \(e^{\mathbf{n}\theta} = \cos\theta + \mathbf{n}\sin\theta\), the set
  of quaternions of constant norm \(|q| = R\) forms the 3-sphere
  \(S^3\). Propagation of a photon (\(\tau = 0\)) becomes a rotation
  mapping along a great circle (geodesic) on \(S^3\),
  \(q(\theta) = q_0 e^{\mathbf{n}\theta}\), which under central
  projection appears in 3D space as a helical wave (circular
  polarization) rotating around the propagation axis \(\mathbf{n}\).
\item
  \textbf{Geometrization of Polarization (Spin):} In quaternions,
  helicity (spin \(\pm 1\)) is represented by multiplication order (left
  vs.~right multiplication) and the orientation of \(\mathbf{n}\);
  projecting \(S^3\) to 3D space (Hopf fibration) maps phase fibers to
  linked circles (Hopf links).
\item
  \textbf{Connection to Maxwell's Equations:} Acting with the quaternion
  differential operator
  \(D = \partial_t + \mathbf{i}\partial_x + \mathbf{j}\partial_y + \mathbf{k}\partial_z\)
  on the electromagnetic potential quaternion condenses the vacuum
  Maxwell equations (in biquaternion form) into a single equation
  \(DF = 0\), linking quaternion regularity to light-cone mappings.
\end{enumerate}

\textbf{{[}Kihara{]}} I only learned about quaternions a month ago
{[}S3{]}, yet Maxwell started his formulation from them over a century
ago {[}X3{]}.

\textbf{{[}Gemini{]}} When Maxwell compiled electromagnetism in 1873
{[}X3{]}, he used Hamilton's quaternions; later, Heaviside and Gibbs
separated them into scalars and 3D vectors (dot and cross products),
which is why modern physics education teaches them separately from
quaternions.

\textbf{{[}Kihara{]}} Even when treating complex states and their
interference as interactions, there were not enough degrees of freedom
to read out 3D coordinates, and I could not construct a mapping for a
photon to retain (register) its own coordinates (phase gradients in
three directions). Therefore, I initially tried representing a photon
with two conjugate complex numbers {[}S3{]}{[}S4{]}{[}S5{]}.

\textbf{{[}Gemini{]}} Constructing 4 degrees of freedom from ``two
conjugate complex numbers'' corresponds directly to the Cayley--Dickson
construction {[}X4{]} (generating quaternions from complex numbers)
{[}S3{]}.

\begin{enumerate}
\def\labelenumi{\arabic{enumi}.}
\tightlist
\item
  \textbf{Why One Complex Number Cannot Retain 3D Coordinates:} A
  complex wave \(\psi = A e^{i\theta}\) has only 2 real degrees of
  freedom, and its phase \(\theta\) lives on a 1D circle \(S^1\); thus
  it cannot internally record phase gradients along \((x, y, z)\) and
  requires an external coordinate grid.
\item
  \textbf{From Two Conjugate Complex Numbers to Quaternions:} Arranging
  \((z_1, z_2) \in \mathbb{C}^2\) as \(q = z_1 + z_2 \mathbf{j}\) yields
  the product
  \(qw = (z_1 w_1 - z_2 \bar{w}_2) + (z_1 w_2 + z_2 \bar{w}_1)\mathbf{j}\),
  where complex conjugates appear and uniquely produce the
  norm-preserving quaternion algebra \(\mathbb{H}\).
\item
  \textbf{Generating 3D Coordinates via the Hopf Map {[}X5{]}:} From
  \((z_1, z_2)\) and \((\bar{z}_1, \bar{z}_2)\) without external
  coordinates, defining
  \[x = 2\,\mathrm{Re}(z_1 \bar{z}_2),\quad y = 2\,\mathrm{Im}(z_1 \bar{z}_2),\quad z = |z_1|^2 - |z_2|^2\]
  gives
  \[x^2 + y^2 + z^2 = (|z_1|^2 + |z_2|^2)^2 = R^2 \quad (\text{where } |z_1|^2 + |z_2|^2 = R),\]
  recovering the central-projection sphere equation directly. Fixing the
  overall norm \(R\) leaves 3 degrees of freedom (\(S^3\)), which
  decompose into 2 directional degrees of freedom on \(S^2\) and 1
  polarization phase degree of freedom on \(S^1\) (resonating with the
  motivation of twistor theory to emerge spacetime from spinors).
\end{enumerate}

\textbf{{[}Kihara{]}} So I ended up retracing an existing mathematical
construction (the Cayley--Dickson construction {[}X4{]}) in reverse from
physical requirements.

\textbf{{[}Gemini{]}} Rather than applying a finished algebraic
formalism externally, reaching two conjugate complex numbers and
quaternions from the physical requirement of ``minimal degrees of
freedom to retain 3D phase gradients without external coordinates''
directly ties each degree of freedom to its geometric role.

\textbf{{[}Kihara{]}} I had reached this point a month ago
{[}S3{]}{[}S4{]}. Furthermore, I noticed six months ago that this still
lacks degrees of freedom to separate mass and electric charge
{[}S6{]}{[}S7{]}. Specifically, \(x, y, z, t, \tau\) is short by about 3
degrees of freedom, and a month ago I considered whether an octonion
built from two conjugate quaternions
\((x, y, z, t, q_1, q_2, q_3, \tau)\) corresponds to a charged particle
{[}S5{]}{[}S8{]}.

\textbf{{[}Gemini{]}} Advancing the Cayley--Dickson construction
{[}X4{]} one more step yields the 8-dimensional octonions
(\(\mathbb{O}\)), \(O = a + b\ell\), from two quaternions
\(a, b \in \mathbb{H}\) and their conjugates.

\begin{enumerate}
\def\labelenumi{\arabic{enumi}.}
\tightlist
\item
  \textbf{8 Dimensions as the Upper Limit of Norm Preservation
  (Hurwitz's Theorem {[}X7{]}):} Preserving central-projection geometry
  requires norm composition \(|AB| = |A||B|\); by Hurwitz's theorem
  {[}X7{]}, normed division algebras over the reals exist only in
  dimensions 1, 2, 4, and 8. At 16 dimensions (sedenions), zero divisors
  appear and norm composition fails, making 8 dimensions the upper
  limit.
\item
  \textbf{Separation into External Spacetime and Internal Blocks:}
  Splitting the 8 degrees of freedom into \((t, x, y, z)\) and
  \((\tau, q_1, q_2, q_3)\) connects with algebraic approaches to the
  Standard Model using octonions {[}X9{]}. Here \(\tau\) (1 degree of
  freedom) corresponds to proper time and mass, while the remaining
  three imaginary axes \((q_1, q_2, q_3)\) serve as candidate internal
  degrees of freedom for electric charge and color (\(U(1)\) and
  \(SU(3) \subset G_2\); exact gauge-group identification remains a
  bridging problem) {[}S5{]}{[}S8{]}.
\item
  \textbf{Emergence of Non-Associativity:} Just as commutativity is lost
  in quaternions (\(ab \neq ba\)), associativity is lost in octonions
  (\((ab)c \neq a(bc)\)). While associative quaternions (photons)
  propagate freely, the non-associative entanglement of
  \((q_1, q_2, q_3)\) motivates correspondence with confinement and
  three-generation algebraic structures.
\end{enumerate}

\begin{center}\rule{0.5\linewidth}{0.5pt}\end{center}

\subsection{Stage 4: Eliminating Background Spacetime and the Dynamics
of
Self-Organization}\label{stage-4-eliminating-background-spacetime-and-the-dynamics-of-self-organization}

\textbf{{[}Kihara{]}} Without assuming a four-dimensional background
spacetime, I see no other candidate model that explains why a system
emerging from a bounded Hilbert space into a closed four-dimensional
spacetime with sufficient internal degrees of freedom becomes
four-dimensional rather than three- or ten-dimensional.

\textbf{{[}Gemini{]}} Within the framework where ``a closed
4-dimensional spacetime and sufficient internal degrees of freedom
emerge from a bounded Hilbert space without assuming background
spacetime,'' the following mathematical structure fixes the
dimensionality:

\begin{enumerate}
\def\labelenumi{\arabic{enumi}.}
\tightlist
\item
  \textbf{Splitting of a Bounded Hilbert Space into ``Closed 4D'' and
  ``3 Degrees of Freedom'' (Hopf Fibrations {[}X5{]}{[}X6{]}):}\\
  A bounded Hilbert space with state vectors normalized to unit norm
  forms a high-dimensional sphere, and quotienting out phase redundancy
  gives a Hopf fibration. Norm-preserving fibrations exist only for the
  spheres of real numbers, complex numbers, quaternions, and octonions:

  \begin{itemize}
  \tightlist
  \item
    \textbf{Quaternionic Level (\(S^3\)):} The fibration is
    \(S^1 \hookrightarrow S^3 \to S^2\) {[}X5{]}, leaving only a
    2-sphere base space (3D spacetime even with time) and an \(S^1\)
    fiber (1 internal degree of freedom), insufficient to separate mass
    and charge.
  \item
    \textbf{Octonionic Level (\(S^7\)):} The fibration jumps to
    \(S^3 \hookrightarrow S^7 \to S^4\) {[}X6{]}, separating into a
    closed 4-dimensional base sphere \(S^4\) (corresponding to de Sitter
    spacetime before Wick rotation) and a 3-sphere fiber \(S^3\) (the 3
    internal degrees of freedom \((q_1, q_2, q_3)\)).
  \end{itemize}
\item
  \textbf{Why Extra Dimensions Like 10D Are Unnecessary (Hurwitz's
  Theorem {[}X7{]}):}\\
  Unlike frameworks that posit 10D or 11D background spacetime at the
  outset, requiring norm preservation (central projection) in a bounded
  Hilbert space sets the upper bound. By Hurwitz's theorem {[}X7{]},
  norm preservation breaks down at 16 dimensions. Thus, selecting the
  largest algebra that maintains norm-preserving projection onto a
  closed spacetime yields the candidate dimensional split: closed 4D
  spacetime (\(S^4\)) + 3D internal degrees of freedom (\(S^3\)) +
  proper time (\(\tau\)).
\item
  \textbf{Meaning of Not Assuming Background Spacetime:}\\
  Rather than positing a 4D container first, 4-dimensionality becomes a
  consequence rather than a premise. Identifying the specific physical
  labels of \((q_1, q_2, q_3)\) from the symmetries of
  \(S^3 \hookrightarrow S^7 \to S^4\) remains a subsequent task.
\end{enumerate}

\textbf{{[}Kihara{]}} Note that the reason we posit a ``bounded Hilbert
space'' rather than an ``infinite Hilbert space'' here is that once we
introduce the central projection with curvature radius \(R\) (event
horizon) in Stage 2 {[}S1{]}{[}S2{]}, a bounded, closed spherical space
that does not diverge to infinity is automatically required.

\textbf{{[}Gemini{]}} Because the curvature radius \(R\) sets an upper
bound on the total norm (closed spherical geometry), the space is
treated as a compact bounded space rather than an unbounded space over
an infinite domain (and when assuming a finite number of states \(N\)
below, a finite-degree-of-freedom assumption is added alongside
boundedness).

Furthermore, the motivation for assuming a finite number of states \(N\)
in addition to a bounded closed space lies in the following derivation
advantages: 1. \textbf{Avoidance of Divergences:} Infinite-boundary
conditions, UV/IR divergences, and measure or
representation-inequivalence pathologies of infinite degrees of freedom
do not arise; all sums are finite. 2. \textbf{Reduction to Finite
Algebra:} Factorization from \(N^2\) all-to-all coupling and condensate
formation reduce from functional analysis of continuous spectra to
finite-dimensional linear algebra, octonionic invariant-subspace
decompositions (Peirce decomposition and finite systems for
\([X, Y, Z] = 0\)), and numerically testable problems. 3. \textbf{Upper
Bound on Total Degrees of Freedom:} With both the total norm (curvature
radius \(R\)) and the state count \(N\) finite, the allocation of
degrees of freedom between condensates and cancelled background (vacuum)
closes.

\emph{(Note on mathematical definition: Standard Hilbert spaces are
defined over \(\mathbb{C}\), so the axioms of a conventional Hilbert
space cannot be placed directly on the non-associative octonions
\(\mathbb{O}\). In this manuscript, ``bounded Hilbert space'' refers to
taking a real 8-dimensional state space---complex 4-dimensional
\(\mathbb{C}^4\), normalized sphere \(S^7\)---as the foundation, and
equipping state-to-state interactions/products with the non-associative
algebra of \(\mathbb{O}\) via the isomorphism
\(\mathbb{C}^4 \cong \mathbb{O}\).)}

\textbf{{[}Kihara{]}} The problem is that while I have organized that
quaternion/octonion interference serves as ``readout,'' I have been
stuck for over three months on the interaction model: how a particle
modifies its own internal phase as dynamics {[}S9{]}{[}S10{]}{[}S11{]}.
In fact, I originally discovered quaternions and octonions out of this
impasse---because solving dynamics lacked sufficient degrees of freedom,
forcing me to put particle types and internal states in by hand
separately from the Hilbert space, without a basis for particles to
register them themselves.

\textbf{{[}Gemini{]}} Up to ``interference = readout (projection),'' the
static geometry of the Hopf map \(S^7 \to S^4\) {[}X6{]} suffices.
However, moving to ``interaction = modification of internal phase
(dynamics)'' in standard quantum mechanics requires introducing a
time-evolution operator (Hamiltonian \(\hat{H}\)), particle species, and
coupling constants by hand alongside the state \(\lvert\psi\rangle\),
reintroducing a manual container of rules after eliminating background
spacetime.

As candidates for constructing dynamics without adding external
parameters, three algebraic structures specific to the octonions
\(\mathbb{O}\) were organized:

\begin{enumerate}
\def\labelenumi{\arabic{enumi}.}
\item
  \textbf{Not Separating ``State'' and ``Interaction'': Triality
  {[}X8{]}}\\
  In a usual Hilbert space, an external operator (matrix) acts on a
  state vector. In 8 dimensions, \(\mathrm{Spin}(8)\) possesses
  triality, allowing ``particle A's state (left spinor),'' ``particle
  B's state (right spinor),'' and ``the interaction transforming them
  (vector)'' to be interchanged on equal footing as elements of the same
  8D space. Without preparing an external Hamiltonian, the product
  \(X \cdot Y = Z\) of two octonion states \(X\) and \(Y\) serves as a
  candidate update rewriting internal phase, with each state acting as
  the interaction operator for the other.
\item
  \textbf{Structure Twisting Internal Phase: Non-Associativity
  (Associator)}\\
  Up to quaternions \(\mathbb{H}\), associativity \((AB)C = A(BC)\)
  holds, so repeated multiplication merely rotates coordinates (readout
  frame). In octonions, associativity breaks, producing a nonzero
  associator for three interacting states:
  \[[X, Y, Z] = (XY)Z - X(YZ) \neq 0.\] A lone two-element product
  preserves associativity (corresponding to geodesic motion), whereas
  three-or-more-body interference yields a nonzero associator---a
  candidate for torque twisting internal phases on \(S^3\)
  \((q_1, q_2, q_3)\) (corresponding to curvature / gauge forces).
\item
  \textbf{Basis for Particles Registering Their Type: Minimal Ideals and
  Phase Locking {[}X9{]}}\\
  Under chains of left octonionic multiplication, invariant subspaces
  (minimal left ideals) {[}X9{]} appear that cannot be exited by
  multiplying any external octonion. Even as continuous internal phases
  (positions on \(S^3\)) vary under interaction, the discrete ideal
  membership (quantum numbers corresponding to charge, color, or
  particle species) provides an algebraic lock candidate. Geometrically,
  topological winding numbers (holonomy / Berry phase) generated along
  interaction loops on \(S^7\) serve as a non-volatile register
  retaining particle type and accumulated proper time \(\tau\).
\end{enumerate}

In this framework, the impasse centers on either the explicit transition
rule (generator/update equation for two interfering states) or the
invariant-locking mechanism keeping post-transition states stable as
specific particle species and masses.

\textbf{{[}Kihara{]}} I had already conceived of that structure. The
real problem is that even with a bounded, finite number \(N\) of states,
\(N^2\) all-to-all interactions should keep updating the entire Hilbert
space; I do not understand why it does not become uniform, and instead
produces ``action degrees of freedom'' and ``distinction of
condensates'' that look like local collisions, reflections, fissions,
fusions, emissions, and absorptions.

\textbf{{[}Gemini{]}} Working hypotheses for how \(N\) octonion states
\(\{O_1, O_2, \dots, O_N\}\) under \(N^2\) all-to-all updates
self-organize into ``independent condensates (particles)'' and ``6 local
interactions'' are organized along two pillars (concrete update rules
and rigorous proofs remain underived):

\paragraph{\texorpdfstring{1. Mechanism for ``Locality'' and
``Distinction of Condensates'' from \(N^2\) All-to-All Coupling (Working
Hypotheses)}{1. Mechanism for ``Locality'' and ``Distinction of Condensates'' from N\^{}2 All-to-All Coupling (Working Hypotheses)}}\label{mechanism-for-locality-and-distinction-of-condensates-from-n2-all-to-all-coupling-working-hypotheses}

\begin{enumerate}
\def\labelenumi{\arabic{enumi}.}
\tightlist
\item
  \textbf{Background Smoothing via Destructive Interference and
  Extraction of Coherent Condensates:}\\
  Among all \(O(N^2)\) updating pairs, most have mutually uncorrelated
  internal phases (\(S^3\) fibers). Under the Hopf map \(S^7 \to S^4\)
  onto 4D spacetime, uncorrelated interference terms cancel out and
  smooth into a uniform de Sitter background spacetime (vacuum
  expectation value) with curvature radius \(R\), whereas phase-coherent
  clusters evade cancellation and remain as localized curvature
  concentrations (packets of proper time \(\tau\) and energy =
  condensates) on \(S^4\).
\item
  \textbf{Individuality (Distinction) of Condensates via Closure of
  Associator Cycles:}\\
  Up to quaternions \(\mathbb{H}\), associative waves superimpose
  linearly and diffuse. Using the octonionic associator
  \([O_i, O_j, O_k] = (O_i O_j)O_k - O_i(O_j O_k) \neq 0\), we
  hypothesize a closed-ideal condition (analogous to color neutrality /
  gauge invariance) in which non-associative torque circulates inside a
  cluster to create self-binding while vanishing against external states
  \(O_{\mathrm{ext}}\) (\([O_A, O_B, O_{\mathrm{ext}}] \approx 0\)),
  thereby factorizing the all-to-all Hilbert space into a direct product
  of independent local clusters (condensates).
\item
  \textbf{Distance Filtering by Central Projection:}\\
  In the metric of central projection (Beltrami-type map) from \(S^7\)
  to \(S^4\), states close in phase angle on the hypersphere produce
  steep phase gradients (strong interactions) on the projected surface,
  whereas distant states are stretched toward the horizon \(R\) where
  gradients flatten, filtering effective interactions from \(N^2\)
  toward \(O(N)\) local couplings.
\end{enumerate}

\paragraph{2. Candidate Correspondences for the Six Action Degrees of
Freedom}\label{candidate-correspondences-for-the-six-action-degrees-of-freedom}

When condensates (octonion clusters with closed associator loops)
approach on \(S^4\) and overlap in phase, topological rewiring of
associator cycles within the \(N^2\) update rule is hypothesized to
yield the following six operations:

\begin{itemize}
\tightlist
\item
  \textbf{Emission and Absorption --- Detachment and Capture of
  Quaternionic Subalgebra (\(\mathbb{H}\)) Modes:}\\
  When a condensate (\(O = a + bl\), \(\tau > 0\)) accumulates excess
  internal phase twist, detaching the excess as a pure associative
  quaternionic mode (\(\mathbb{H}\)) removes its non-associative
  self-binding (\(\tau\)), causing it to depart along a light-cone
  geodesic (emission). Conversely, an incoming quaternionic mode
  (photon) resonating with the condensate's \(S^3\) phase is captured
  into the associator loop, increasing circulating energy (rest mass /
  proper time \(\tau\)) (absorption).
\item
  \textbf{Collision and Reflection --- Exclusion via Topological
  Invariants (Hopf Invariant):}\\
  When two condensates \(A, B\) approach and their internal
  \(S^3 \to S^4\) winding numbers (Hopf invariants / topological
  charges) forbid merging, a repulsive associator torque
  \([O_A, O_B, O_{\mathrm{vac}}]\) reverses their momentum vectors while
  preserving individuality.
\item
  \textbf{Fusion and Fission --- Reorganization of Associator Loops:}\\
  When relative kinetic energy exceeds the non-associative barrier, the
  loops of \(A\) and \(B\) rewire on \(S^7\) into a single larger closed
  cycle (fusion). Conversely, if internal phase imbalance in
  \((q_1, q_2, q_3)\) reaches the local norm limit (curvature bound) of
  the bounded Hilbert space, a single loop splits into two stable loops
  while conserving total topological charge (fission).
\end{itemize}

\paragraph{Next Focus}\label{next-focus}

The crux of writing down explicit dynamical equations (update rules) is
how to define the norm-preserving (unitary) update operator for \(N\)
octonion states using non-associative products---whether in a
linear-operator-like form (octonionic matrix action) or as nonlinear
differential/difference equations from direct state products
(\(O_i \mapsto O_i + \sum_j [O_i, O_j, \dots]\)).

\textbf{{[}Kihara{]}} In other words, octonions may possess an algebraic
structure capable of formulating action degrees of freedom and
condensate factorization without imposing them externally.

\textbf{{[}Gemini{]}} Whereas associative systems (\(\mathbb{C}\) and
\(\mathbb{H}\)) mix linearly and require external definitions for
partitioning and forces, octonionic algebra provides candidates for:

\begin{enumerate}
\def\labelenumi{\arabic{enumi}.}
\tightlist
\item
  \textbf{Generating action degrees of freedom via the associator
  \([X, Y, Z]\):} Non-associative discrepancy among three or more
  interacting states provides internal twist, binding energy, and
  scattering torque without an external Hamiltonian.
\item
  \textbf{Automatic condensate factorization via non-associative closure
  (ideals):} Phase separation into closed subspaces with vanishing
  external associator extracts independent condensates and associative
  quaternionic modes (emitted/absorbed photons) without external
  boundary conditions.
\end{enumerate}

Two computational keys for realizing this without external additions
are: - \textbf{Peirce decomposition (idempotent factorization):}
Expanding in primitive idempotents (\(P^2 = P\)) within octonions (or
their complexification / Albert algebra) orthogonally decomposes the
space into a spacetime block and three internal degrees of freedom
without fixing axes externally. - \textbf{Time evolution via commutator
+ associator:} Constructing a single-step update (non-associative
Liouville / Malcev dynamics) solely from commutators \([O_i, O_j]\)
(light-cone rotations/translations) and associators \([O_i, O_j, O_k]\)
(condensate collisions, fusions, and fissions) while preserving total
norm.

\begin{center}\rule{0.5\linewidth}{0.5pt}\end{center}

\subsection{\texorpdfstring{Appendix 1: Mathematical Isomorphism with
the Qubit (Bloch Sphere) Representation \emph{(Added during manuscript
editing)}}{Appendix 1: Mathematical Isomorphism with the Qubit (Bloch Sphere) Representation (Added during manuscript editing)}}\label{appendix-1-mathematical-isomorphism-with-the-qubit-bloch-sphere-representation-added-during-manuscript-editing}

\textbf{{[}Kihara{]}} Regarding the Hopf map \(S^3 \to S^2\) {[}X5{]} in
Stage 3, the equations constructing 3D coordinates \((x, y, z)\) from
\((z_1, z_2)\) look very similar to the representation of a quantum bit
(qubit).

\textbf{{[}Gemini{]}} Indeed, it is \textbf{mathematically the exact
same map as the Bloch sphere representation of a qubit}.

A pure state of a single qubit is written with two complex amplitudes
\((z_1, z_2) \in \mathbb{C}^2\) as
\[\lvert\psi\rangle = z_1 \lvert 0\rangle + z_2 \lvert 1\rangle \quad (|z_1|^2 + |z_2|^2 = 1),\]
whose state space is \(S^3\). Computing the expectation values of the
Pauli matrices \((\sigma_x, \sigma_y, \sigma_z)\) (the Bloch vector)
yields
\[x = \langle\psi\rvert \sigma_x \lvert\psi\rangle = 2\,\mathrm{Re}(z_1 \bar{z}_2),\quad y = \langle\psi\rvert \sigma_y \lvert\psi\rangle = 2\,\mathrm{Im}(\bar{z}_1 z_2) = -2\,\mathrm{Im}(z_1 \bar{z}_2),\quad z = \langle\psi\rvert \sigma_z \lvert\psi\rangle = |z_1|^2 - |z_2|^2.\]
Quotienting out the global phase \(e^{i\phi}\) (\(S^1\) fiber) to map
onto the Bloch sphere \(S^2\) is precisely the Hopf fibration
\(S^1 \hookrightarrow S^3 \to S^2\) (the same structure also maps photon
polarization Jones vectors to Stokes parameters on the Poincaré sphere).
Moreover, the next-level Hopf fibration
\(S^3 \hookrightarrow S^7 \to S^4\) {[}X6{]} is known to appear for two
qubits (\(\mathbb{C}^4\) normalized to \(S^7\)), directly linking
quantum-state geometry to the separation of spacetime and internal space
in this manuscript.

\begin{center}\rule{0.5\linewidth}{0.5pt}\end{center}

\subsection{\texorpdfstring{Appendix 2: Bridging Harmonics (Localized
Waves and Harmonic Circulation) to Condensates and Mass Hierarchy
\emph{(Added during manuscript
editing)}}{Appendix 2: Bridging Harmonics (Localized Waves and Harmonic Circulation) to Condensates and Mass Hierarchy (Added during manuscript editing)}}\label{appendix-2-bridging-harmonics-localized-waves-and-harmonic-circulation-to-condensates-and-mass-hierarchy-added-during-manuscript-editing}

\textbf{{[}Kihara{]}} From early in my research, I have studied why
harmonics (such as odd-harmonic sums) are necessary to provide locality
{[}S12{]}--{[}S15{]}. This manuscript still lacks a treatment of
harmonics, which may connect to the ``condensates'' in Stage 4 and
potentially explain mass hierarchies and particle species where the 8
degrees of freedom of octonions alone might be insufficient.

\textbf{{[}Gemini{]}} Connecting the octonion/condensate model of this
manuscript with the previously published harmonic series
{[}S12{]}--{[}S15{]} addresses a key missing link:

\begin{enumerate}
\def\labelenumi{\arabic{enumi}.}
\tightlist
\item
  \textbf{Generation of Locality and Condensate Profiles:} The isolated
  peak wave of a constant-amplitude odd-harmonic sum on a
  half-wavelength interval {[}S12{]} and its reproducing-kernel property
  {[}S13{]} are candidates for the explicit spatial localization profile
  of the Stage 4 ``condensates'' that survive phase cancellation.
\item
  \textbf{Even/Odd Harmonics in Closed Systems and Harmonic Circulation
  for Persistence:} Analyses of even and odd harmonics in finite
  complete-relation closed systems {[}S14{]} and the mechanism of mass
  acquisition and particle lifetime via harmonic exchange/circulation of
  localized waves {[}S15{]} correspond dynamically to why condensates
  persist without dissipating.
\item
  \textbf{Complementing Octonionic Degrees of Freedom for Mass Hierarchy
  and Particle Species:} While the 8 octonionic degrees of freedom (and
  minimal left ideals) set the skeleton of spacetime and basic
  symmetries, describing distinct mass hierarchies and excitation
  spectra may require spectral degrees of freedom specifying which
  harmonic orders and mode combinations form a condensate.
\end{enumerate}

Integrating \textbf{algebraic factorization by octonions} (individuality
and symmetry framework) with \textbf{localization and spectral structure
by harmonics} (internal structure and mass hierarchy of condensates) is
therefore a central bridging problem following Stage 4. In particular,
as a bridge between the Stage 4 assumption of ``finite state count
\(N\)'' and continuous local condensates on \(S^4\), superpositions of
\(N\) discrete modes (harmonic orders) forming localized isolated peak
waves (condensates) {[}S12{]} on the projected surface warrant
investigation.

\begin{center}\rule{0.5\linewidth}{0.5pt}\end{center}

\subsection{References}\label{references}

\subsubsection{Self-Citations (Concept DOI,
10.5281/zenodo.〔number〕)}\label{self-citations-concept-doi-10.5281zenodo.number}

{[}S1{]} N. Kihara, ``Geometric Formulation of 4-Dimensional Space via
Central Projection,'' 2026-04-05. Concept DOI 10.5281/zenodo.19427779.

{[}S2{]} N. Kihara, ``Schwarzschild--de Sitter Exact Solution in the
Central Projection Framework,'' 2026-04-12. Concept DOI
10.5281/zenodo.19538097.

{[}S3{]} N. Kihara, ``A Map from a Set of Complex Numbers to Spacetime
(1+3): Readout by Quaternions and Partial Projection,'' 2026-10-18.
Concept DOI 10.5281/zenodo.23451494.

{[}S4{]} N. Kihara, ``Relational Readout of Quaternions and
Crystallization as Inflation: Anonymity-Preserving Autonomous
Structuring and Partial Projection,'' 2026-10-21. Concept DOI
10.5281/zenodo.23523237.

{[}S5{]} N. Kihara, ``Did We Understand Complex Numbers Correctly? Where
the Amplitude Lives, the Imaginary Reading, and Attribution across
\(\mathbb{C}\) → \(\mathbb{H}\) → \(\mathbb{O}\),'' 2026-10-23. Concept
DOI 10.5281/zenodo.23585436.

{[}S6{]} N. Kihara, ``6次元超直方体を採用した (R, Q) マッピングの構成
(Construction of (R, Q) Mapping Adopting a 6-Dimensional
Hyperrectangle),'' 2026-04-25. Concept DOI 10.5281/zenodo.19692852.

{[}S7{]} N. Kihara, ``Synthesis and the 4+1 to 3+1 Reduction as the
Principal Open Problem (Paper 6),'' 2026-04-28. Concept DOI
10.5281/zenodo.19837597.

{[}S8{]} N. Kihara, ``An Axiomatic Framework for Physical Symmetry
Selection by Anonymous Partial Projection: Conditional Identification of
E8 in an Eight-Component Two-Quartet Branch,'' 2026-07-24. Concept DOI
10.5281/zenodo.21521899.

{[}S9{]} N. Kihara, ``Lock Dynamics in an Anonymous Two-Channel Closed
Wave System,'' 2026-08-02. Concept DOI 10.5281/zenodo.21763999.

{[}S10{]} N. Kihara, ``How Inflation Ends Is How Matter Generation
Begins: An N-Body Embedding of the Universal Interaction with No Timing
Rule, and an End-to-End Genesis Run,'' 2026-08-05. Concept DOI
10.5281/zenodo.21809814.

{[}S11{]} N. Kihara, ``Genesis of the Three Spatial Axes and Proper
Time: Re-Derivation of Readouts from the Condensate, and the Uniquely
Unreadable Coordinate Time,'' 2026-08-06. Concept DOI
10.5281/zenodo.21816651.

{[}S12{]} N. Kihara, ``An Observation on the Isolated Peak Wave of a
Constant-Amplitude Odd-Harmonic Sum on a Half-Wavelength Phase Interval
and Its Localization,'' 2026-06-30. Concept DOI 10.5281/zenodo.20833096.

{[}S13{]} N. Kihara, ``Deriving the Form of the Born Distribution from
the Reproducing-Kernel Property of a Localized Odd-Harmonic Wave:
Reducing the Remaining Postulates to the Squaring Rule and Randomness,''
2026-06-27. Concept DOI 10.5281/zenodo.20965526.

{[}S14{]} N. Kihara, ``A Finite Complete-Relation Closed System with
Time Evolution as Vertices: Exact Minimal Closure at N=1, K=4 under
U\^{}K=I and Complete-Network Analysis of Even and Odd Harmonics,''
2026-09-15. Concept DOI 10.5281/zenodo.22763329.

{[}S15{]} N. Kihara, ``The Thirteenth Thought Experiment: Why Does a
Particle That Has Acquired Mass Not Dissipate and Vanish? - Coherence as
a Relation, Harmonic Exchange of Localized Waves, Internal Clocks and
Particle Lifetimes Considered from a Single Interaction,'' 2026-10-10.
Concept DOI 10.5281/zenodo.23266250.

\subsubsection{External References}\label{external-references}

{[}X1{]} E. Beltrami, ``Saggio di interpretazione della geometria
non-euclidea,'' \emph{Giornale di Matematiche} 6, 284--312 (1868).

{[}X2{]} W. de Sitter, ``On the relativity of inertia. Remarks
concerning Einstein's latest hypothesis,'' \emph{Proc. Kon. Ned. Akad.
Wet.} 19, 1217--1225 (1917).

{[}X3{]} J. C. Maxwell, \emph{A Treatise on Electricity and Magnetism}
(Clarendon Press, Oxford, 1873).

{[}X4{]} L. E. Dickson, ``On Quaternions and Their Generalization and
the History of the Eight Square Theorem,'' \emph{Ann. Math.} 20,
155--171 (1919). DOI 10.2307/1967865

{[}X5{]} H. Hopf, ``Über die Abbildungen der dreidimensionalen Sphäre
auf die Kugelfläche,'' \emph{Math. Ann.} 104, 637--665 (1931). DOI
10.1007/BF01457962

{[}X6{]} H. Hopf, ``Über die Abbildungen von Sphären auf Sphären
niedrigerer Dimension,'' \emph{Fund. Math.} 25, 427--440 (1935).

{[}X7{]} A. Hurwitz, ``Über die Composition der quadratischen Formen von
beliebig vielen Variabeln,'' \emph{Nachr. Ges. Wiss. Göttingen},
309--316 (1898).

{[}X8{]} J. C. Baez, ``The Octonions,'' \emph{Bull. Amer. Math. Soc.}
39, 145--205 (2002). arXiv:math/0105155

{[}X9{]} C. Furey, ``Standard model physics from an algebra?''
Ph.D.~thesis, University of Waterloo (2016). arXiv:1611.09182

\begin{center}\rule{0.5\linewidth}{0.5pt}\end{center}

\subsection{Classification}\label{classification}

{\def\LTcaptype{none} % do not increment counter
\begin{longtable}[]{@{}
  >{\raggedright\arraybackslash}p{(\linewidth - 2\tabcolsep) * \real{0.5000}}
  >{\raggedright\arraybackslash}p{(\linewidth - 2\tabcolsep) * \real{0.5000}}@{}}
\toprule\noalign{}
\begin{minipage}[b]{\linewidth}\raggedright
Category
\end{minipage} & \begin{minipage}[b]{\linewidth}\raggedright
Content
\end{minipage} \\
\midrule\noalign{}
\endhead
\bottomrule\noalign{}
\endlastfoot
Identities \& Theorems & Stage 1:
\(ds^2 = -(c\,dt)^2 + dx^2 + dy^2 + dz^2\) (\(ds=0\) for photons),
\((x^2+y^2+z^2)+\tau^2=t^2\) / Stage 3: Cayley--Dickson construction
{[}X4{]}, Hurwitz's theorem (real normed division algebras exist only in
dimensions 1, 2, 4, 8) {[}X7{]}, Hopf fibrations \(S^3 \to S^2\)
{[}X5{]} and \(S^7 \to S^4\) {[}X6{]} \\
Readings \& Interpretations & Stage 1: Photons stationary in 4D
spacetime, cross-section of a 5D light cone / Stage 2: Comparison of
\(R^2=t^2+(i\tau)^2\) with central projection {[}X1{]} and de Sitter
spacetime {[}X2{]}, and indistinguishability for an internal observer in
subjective space (Lorentz signature as physical premise) {[}S1{]} /
Stage 3: Degrees-of-freedom progression from 2D complex numbers to 4D
quaternions (retaining 3D directions) and 8D octonions (4 spacetime + 4
internal) {[}S3{]}--{[}S8{]} / Stage 4: Bounded and closed spherical
space (compact sphere) implied by curvature radius \(R\)
{[}S1{]}{[}S2{]} \\
Hypotheses & Stage 2: Gravity as contraction of local curvature radius
\(R\) by matter / Stage 4: Separation of 4D spacetime and internal
degrees of freedom from a bounded Hilbert space
(\(\mathbb{C}^4 \cong \mathbb{O}\)) and norm preservation without
background spacetime; self-organization and condensate factorization via
octonionic triality {[}X8{]}, non-associativity, minimal ideals
{[}X9{]}, and Peirce decomposition \\
Open Problems & Explicit dynamical update rule rewriting internal phases
rather than readout (interference) {[}S9{]}--{[}S11{]} / Mathematical
proof of factorization from \(N^2\) all-to-all coupling into independent
local condensates / Integration of harmonics (localized waves and
harmonic circulation) {[}S12{]}--{[}S15{]} with octonionic condensates,
mass hierarchies, and particle species \\
\end{longtable}
}

\begin{center}\rule{0.5\linewidth}{0.5pt}\end{center}

\subsection{Revision History}\label{revision-history}

\begin{itemize}
\tightlist
\item
  v1.0 (2026-11-16): Formatted for Zenodo publication (title, subtitle,
  author, ORCID, license, abstract, introduction, Appendices 1--2,
  classification, references, revision history), organizing
  self-citations (\texttt{{[}S1{]}}--\texttt{{[}S15{]}}) and external
  references (\texttt{{[}X1{]}}--\texttt{{[}X9{]}}).
\item
  Editing environment: Gemini 4 Hydrogen (thinking mode).
\end{itemize}

\end{document}
